\documentclass[10pt, aspectratio=169]{beamer}

% --- Configuración del Tema ---
\usetheme{metropolis}
\usecolortheme{owl} 
\usepackage[spanish]{babel}
\usepackage[utf8]{inputenc}
\usepackage{amsfonts, amsmath, amssymb, amsthm}
\usepackage{graphicx}
\usepackage{tikz}

% --- Colores Personalizados ---
\definecolor{diagonalred}{RGB}{140, 20, 40}

% --- Información del Título ---
\title{Los Límites de la Razón}
\subtitle{Sesión 6: El Colapso de lo Abstracto en lo Físico}
\author{Rosh Guadiana}
\date{Junio 2026}

\begin{document}

\maketitle

% ==========================================
% SHANNON Y EL FINITISMO
% ==========================================

\begin{frame}{El Isomorfismo de Shannon}
    \vspace{0.3cm}
    \Large \textbf{La respuesta finita a Hilbert}
    
    \vspace{0.4cm}
    \normalsize
    El paso del espacio de estados infinito de Turing al espacio discreto finito. El Álgebra de Boole se revela como un Anillo Conmutativo Unitario $\langle \mathbb{Z}_2, \oplus, \cdot \rangle$.
    
    \vspace{0.5cm}
    \alert{\textbf{El mapeo físico:}} \\
    La reducción de la Máquina de Turing a compuertas combinacionales binarias en silicio. La equivalencia estructural entre la disyunción exclusiva y la suma en el anillo $\mathbb{GF}(2)$ se formaliza como:
    
    \vspace{0.4cm}
    \Large
    $$A \oplus B = (A \land \neg B) \lor (\neg A \land B)$$
\end{frame}

% ==========================================
% VON NEUMANN Y MATRICES
% ==========================================

\begin{frame}{Von Neumann y el Constructor Universal}
    \vspace{0.3cm}
    \Large \textbf{De la cinemática a la matriz bidimensional}
    
    \vspace{0.4cm}
    \normalsize
    Para eliminar el cabezal móvil de Turing y lograr computar estructuras autorreplicantes, Von Neumann mapeó la cinta infinita en un plano discreto. 
    
    \vspace{0.5cm}
    \alert{\textbf{Autómatas Celulares de 29 estados:}} \\
    El ``cabezal'' es ahora una perturbación del estado de las celdas circundantes. Definimos rigurosamente la vecindad de Von Neumann extendida en el retículo:
    
    \vspace{0.4cm}
    \Large
    $$N(x,y) = \{(x+i, y+j) \mid |i| + |j| \le 1\}$$
\end{frame}

% ==========================================
% REGLA 30 Y CANTOR
% ==========================================

\begin{frame}{Regla 30: Diagonalización en Topologías Discretas}
    \vspace{0.2cm}
    \Large \textbf{La irreducibilidad computacional}
    
    \vspace{0.4cm}
    \normalsize
    Un sistema localmente determinado que genera un espacio de configuraciones no numerable. No existen ``atajos gödelianos'' analíticos para saltar el tiempo $t$.
    
    \vspace{0.4cm}
    \begin{columns}
        \begin{column}{0.6\textwidth}
            \alert{\textbf{Evolución local combinatoria (ECA):}}
            \large
            $$a_i^{t+1} = a_{i-1}^t \pmod 2 \oplus (a_i^t \lor a_{i+1}^t)$$
        \end{column}
        \begin{column}{0.4\textwidth}
            \small
            \textbf{El Cono de Luz Computacional:} \\
            La alteración de un bit propaga un error con pendiente máxima $\Delta x / \Delta t = 1$, evidenciando caos determinista.
        \end{column}
    \end{columns}
    
    \vspace{0.4cm}
    % Sugerencia: Aquí puedes insertar el gráfico de la Regla 30
    % \begin{center} \includegraphics[height=0.3\textheight]{../../assets/regla30.png} \end{center}
\end{frame}

% ==========================================
% CONWAY Y TURING
% ==========================================

\begin{frame}{Conway y la Turing-Complejidad en $\mathbb{Z}^2$}
    \vspace{0.3cm}
    \Large \textbf{El Entscheidungsproblem en un tablero de Go}
    
    \vspace{0.4cm}
    \normalsize
    Emergencia topológica: las reglas del Juego de la Vida son capaces de simular cualquier Máquina Universal de Turing mediante puertas lógicas formadas por colisiones de \textit{Gliders}.
    
    \vspace{0.6cm}
    \alert{\textbf{El problema de la parada reencarnado:}} \\
    Resulta matemáticamente indecidible crear un algoritmo general que determine si una configuración inicial arbitraria en este retículo se extinguirá o proliferará infinitamente.
\end{frame}

% ==========================================
% SINTAXIS Y SEMÁNTICA
% ==========================================

\begin{frame}{La Perspectiva de Manuel Morsa}
    \begin{center}
        \Large \textbf{El espejismo del observador}
        
        \vspace{0.8cm}
        \huge \textit{``La inteligencia es un efecto óptico\\ del observador.''}
        \vspace{0.8cm}
    \end{center}
    
    \normalsize
    \alert{\textbf{Sintaxis $\neq$ Semántica:}} \\
    Los autómatas celulares son desiertos semánticos. Ejecutan transformaciones homotópicas de cadenas de bits. La comprensión es una proyección desde el exterior del sistema.
    
    \vspace{0.4cm}
    \begin{center}
        \large Mapeo Formal [Sintaxis] $\longrightarrow$ Interpretación [Semántica]
    \end{center}
\end{frame}

% ==========================================
% ASFIXIA TERMODINÁMICA
% ==========================================

\begin{frame}{La Asfixia Termodinámica}
    \vspace{0.2cm}
    \Large \textbf{El choque contra el Principio de Landauer}
    
    \vspace{0.4cm}
    \normalsize
    La máquina de Turing ideal carece de fricción. En la realidad, \textbf{el bit es un grado de libertad termodinámico}.
    
    \vspace{0.4cm}
    \begin{columns}
        \begin{column}{0.48\textwidth}
            \alert{\textbf{Principio de Landauer:}} \\
            \vspace{0.1cm}
            \small Borrar o cambiar información disipa un mínimo innegociable de calor al entorno.
            \Large
            $$\Delta Q \ge k_B T \ln 2$$
        \end{column}
        \begin{column}{0.48\textwidth}
            \alert{\textbf{Sistemas Lineales:}} \\
            \vspace{0.1cm}
            \small En un sistema lineal ($V_{out} = \alpha V_{in} + \xi$), la varianza del ruido de Johnson-Nyquist ($\sigma^2$) diverge tras $n$ iteraciones de la regla.
        \end{column}
    \end{columns}
    
    \vspace{0.6cm}
    \begin{center}
        \textbf{\color{diagonalred} La entropía física destruye la consistencia sintáctica.}
    \end{center}
\end{frame}

% ==========================================
% NO-LINEALIDAD Y ATRACTORES
% ==========================================

\begin{frame}{La Ley Innegociable de la No-Linealidad}
    \vspace{0.3cm}
    \Large \textbf{El redentor físico del bit}
    
    \vspace{0.4cm}
    \normalsize
    Para que las verdades lógico-matemáticas sobrevivan, el hardware debe violar la proporcionalidad lineal y aplastar las fluctuaciones térmicas activamente.
    
    \vspace{0.4cm}
    \alert{\textbf{Contracción del espacio de fases:}} \\
    Requerimos funciones de transferencia no lineales con regiones asintóticas de saturación donde la ganancia diferencial disipe el ruido hacia atractores fijos ($V_{DD}$ o Tierra):
    
    \vspace{0.3cm}
    \Large
    $$A_v = \frac{dV_{out}}{dV_{in}}, \quad \text{con zonas estables donde } |A_v| < 1$$
\end{frame}

\end{document}
